\section{Introduction}\label{sec:introduction}

Clinical data are multimodal: a chest radiograph travels with its report, a cardiology study with
its discharge letter. Releasing either half requires both to be protected, since an image may need
defacing or removal of burned-in text while its report needs identifier replacement. Artificial
intelligence scales both tasks, and its errors vary by distribution, language, script and identifier
class. Packh\"auser et al.~\cite{packhauser2022reid} re-identified patients from de-identified chest
radiographs and priced a defence against diagnostic utility~\cite{packhauser2023anonymization};
pathological speech has been evaluated against privacy, utility and
fairness~\cite{arasteh2024speaker}; clinical text is instead kept with the model inside the
hospital~\cite{hou2025llama}. We ask what survives transformation of the text, and what protecting
it costs.

\emph{De-identification} reduces identifying information. Reversible \emph{pseudonymisation}
replaces an identifier with a realistic \emph{surrogate}; \emph{anonymisation} prevents protects individuals. 
Stable pseudonyms preserve longitudinal utility but hand an adversary a persistent handle.
European guidance regards keyed message authentication codes as
robust~\cite{enisa2021pseudonymisation}, and the health-informatics standard is revised
accordingly~\cite{din25237}.

Detection benchmarks ask whether an identifying span was
found~\cite{vats2026redact,jha2026piibench,uppala2026privacyfilter}; utility studies measure
downstream change~\cite{vakili2022utility,berg2020impact} or information
loss~\cite{trienes2024infolossqa}; re-identification studies measure attack success. Frameworks
combine technical performance, information loss and intrusion tests~\cite{mozes2021nointruder},
while surrogate systems and masking benchmarks cover single
pieces~\cite{eder2019emails,pilan2022tab,osborne2022bratsynthetic}. No study states together what one
policy buys and what it costs. Data-protection officers release documents and cases, not average
tokens, as one unchanged mention can identify someone after thousands were removed, so we measure
surviving findings, exposed cases and linkable entities rather than token rates alone.

Our design covers four corpora in German, English, Chinese and Arabic: clinical letters, legal
judgments, news and other genres, and e-mail. Every document receives the same detector comparison,
release conditions and measurements. Two predictions are under investigation: that detector precision
governs utility more strongly than the choice between surrogates and typed placeholders, and that
detection sensitivity moves the two attacks in opposite directions, because removing a name obstructs
context linkage while leaving a larger, more stable population for frequency matching. 
